\documentclass[11pt,a4paper]{article}

\usepackage[utf8]{inputenc}
\usepackage[T1]{fontenc}
\usepackage{amsmath,amssymb,amsfonts,amsthm}
\usepackage{mathtools}
\usepackage{booktabs}
\usepackage{array}
\usepackage[table]{xcolor}
\usepackage{geometry}
\usepackage{hyperref}
\usepackage{fancyhdr}
\usepackage{enumitem}
\usepackage{float}
\usepackage{caption}
\usepackage{setspace}
\usepackage{authblk}
\usepackage{algorithm}
\usepackage{algorithmic}

\geometry{margin=2.5cm, top=3cm, bottom=3cm}
\onehalfspacing

\definecolor{headerblue}{RGB}{41,65,122}
\definecolor{lightgray}{RGB}{245,245,245}

\hypersetup{colorlinks=true, linkcolor=headerblue, citecolor=headerblue, urlcolor=headerblue}

\newcolumntype{L}[1]{>{\raggedright\arraybackslash}p{#1}}
\newcolumntype{C}[1]{>{\centering\arraybackslash}p{#1}}

\newtheorem{theorem}{Theorem}
\newtheorem{proposition}[theorem]{Proposition}
\newtheorem{corollary}[theorem]{Corollary}
\newtheorem{lemma}[theorem]{Lemma}
\newtheorem{definition}[theorem]{Definition}
\newtheorem{remark}[theorem]{Remark}

\pagestyle{fancy}
\fancyhf{}
\fancyhead[L]{\small\textit{Unified Homotopy Framework for RBF Identification and ODE Simulation}}
\fancyhead[R]{\small\textit{Rodrigo, Pati\~no}}
\fancyfoot[C]{\thepage}

%% ======================================================================
\title{\textbf{A Unified Homotopy Framework for Nonlinear System Identification and Simulation via Radial Basis Function Regressors}}

\author[1]{Rodolfo H.\ Rodrigo\thanks{Corresponding author. E-mail: rrodrigo@inaut.unsj.edu.ar}}
\author[1]{H.\ Daniel Pati\~no}
\affil[1]{INAUT, Universidad Nacional de San Juan -- CONICET, San Juan, Argentina}

\date{\small Draft --- April 2026\\[4pt]
\small\textit{Target: Neurocomputing / Neural Networks}}

\begin{document}
\maketitle
\thispagestyle{fancy}

%% ======================================================================
\begin{abstract}
We show that Liao's homotopy deformation equation provides a unified algorithmic framework for three problems that are traditionally solved with separate methods: (i)~identifying the parameters of a Radial Basis Function (RBF) network from data, (ii)~solving the resulting nonlinear ODE in real time, and~(iii)~refining the RBF model online as new data arrive. The key observation is that the zeroth-order deformation equation $(1-q)\,\mathcal{L}[\phi - \phi_0] = q\,c_0\,\mathcal{N}[\phi]$ is \emph{universal}: it makes no assumption about the nature of~$\mathcal{N}$. By choosing $\mathcal{L}$ and~$\mathcal{N}$ appropriately, the same homotopy corrections $z_1$, $z_2$, $z_3$---parametrized by~$\gamma = c_0 g_0/\lambda$ and $\sigma = c_0 g'_0/\lambda$---solve all three problems. For identification, $\mathcal{N}$ is the RBF approximation error and $\mathcal{L} = I$ (identity); the linear weights are obtained in a single homotopy step (exact), while the nonlinear parameters (centers, widths) are refined via Newton--Halley corrections. For simulation, $\mathcal{N}$ is the ODE residual $\dot{y} + \mathrm{RBF}(y) - u(t)$ and $\mathcal{L} = d/dt + k$ (linear part of the RBF); the solution departs from the known linear response $y_0(t) = y(0)e^{-kt} + \int_0^t e^{-k(t-s)}u(s)\,ds$ and is corrected step by step. The separation between linear and nonlinear parameters exploits the bilinear structure of the RBF, avoiding the need to optimize all~$3M$ parameters simultaneously. The resulting algorithm uses closed-form corrections at every level, requires no gradient computation, no learning rate, and no backpropagation.
\end{abstract}

\noindent\textbf{Keywords:} Homotopy Analysis Method, radial basis functions, nonlinear system identification, ODE solvers, Newton--Halley corrections.

%% ======================================================================
\section{Introduction}
\label{sec:intro}
%% ======================================================================

Consider a nonlinear dynamical system observed through input--output data:
\begin{equation}
\dot{y}(t) + f(y(t)) = u(t), \qquad y(0) = y_0,
\label{eq:ode}
\end{equation}
where $f$ is an unknown nonlinearity and $u(t)$ is a known excitation. Two classical problems arise: \emph{identification}---find $f$ from data $\{(t_i, y_i, u_i)\}$---and \emph{simulation}---given $f$, compute $y(t)$ for new inputs. In practice, a third problem appears: \emph{online refinement}---update the model of~$f$ as new data arrive during operation.

These three problems are traditionally attacked with separate tools: least squares or backpropagation for identification~\cite{Haykin2009,Park1991}, Runge--Kutta or BDF methods for simulation~\cite{Hairer1993}, and recursive least squares or extended Kalman filters for online update~\cite{Ljung1999}. Each tool has its own convergence theory, parameters, and failure modes.

\textbf{Contribution.} We show that all three problems are instances of a single equation---Liao's zeroth-order deformation equation from the Homotopy Analysis Method (HAM)~\cite{Liao2003,Liao2012}---and that the same closed-form corrections $z_1$, $z_2$, $z_3$ solve all of them. The nonlinearity~$f$ is represented by a Radial Basis Function (RBF) network:
\begin{equation}
\hat{f}(y) = \sum_{j=1}^{M} w_j\,\varphi\!\left(\frac{\|y - c_j\|}{\lambda_j}\right),
\label{eq:rbf}
\end{equation}
where $\mathbf{w} = (w_1,\ldots,w_M)^T$ are linear weights, $\mathbf{c} = (c_1,\ldots,c_M)^T$ are centers, and $\boldsymbol{\lambda} = (\lambda_1,\ldots,\lambda_M)^T$ are widths. The parameter vector is $\boldsymbol{\theta} = (\mathbf{w}, \mathbf{c}, \boldsymbol{\lambda}) \in \mathbb{R}^{3M}$.

The central observation is:

\begin{quote}
\emph{The HAM deformation equation is universal. It solves $\mathcal{N}[\phi] = 0$ for any nonlinear operator~$\mathcal{N}$, by deforming from a known solution~$\phi_0$ of a simpler problem $\mathcal{L}[\phi] = 0$. The nature of~$\phi$ (a function, a vector of parameters, or a trajectory) and the nature of~$\mathcal{N}$ (an approximation error, an ODE residual, or a tracking error) are irrelevant to the algebraic structure of the corrections.}
\end{quote}


%% ======================================================================
\section{The universal deformation equation}
\label{sec:universal}
%% ======================================================================

\subsection{Liao's zeroth-order deformation}

Given a nonlinear operator $\mathcal{N}$ acting on an unknown~$\phi$, the HAM constructs a family $\phi(q)$, $q \in [0,1]$, via:
\begin{equation}
(1-q)\,\mathcal{L}[\phi(q) - \phi_0] = q\,c_0\,\mathcal{N}[\phi(q)],
\label{eq:deformation}
\end{equation}
where $\mathcal{L}$ is a chosen linear operator, $\phi_0$ is an initial approximation, and $c_0 \neq 0$ is a convergence-control parameter. At $q = 0$: $\phi = \phi_0$. At $q = 1$: $\mathcal{N}[\phi] = 0$ (the target).

\subsection{The universal corrections}

Expanding $\phi(q) = \phi_0 + \sum_{m=1}^{\infty} z_m\,q^m$ and Taylor-expanding $\mathcal{N}[\phi_0 + \delta]$, matching powers of~$q$ yields:
\begin{equation}
\boxed{z_1 = \gamma, \qquad z_2 = \gamma(1+\sigma), \qquad z_3 = \gamma(1+\sigma)^2 + \frac{c_0\,\mathcal{N}''_0}{2\lambda}\,\gamma^2,}
\label{eq:z123}
\end{equation}
where
\begin{equation}
\gamma \triangleq \frac{c_0\,\mathcal{N}_0}{\lambda}, \qquad \sigma \triangleq \frac{c_0\,\mathcal{N}'_0}{\lambda},
\label{eq:gamma-sigma}
\end{equation}
with $\mathcal{N}_0 = \mathcal{N}[\phi_0]$, $\mathcal{N}'_0 = \partial\mathcal{N}/\partial\phi\big|_{\phi_0}$, $\mathcal{N}''_0 = \partial^2\mathcal{N}/\partial\phi^2\big|_{\phi_0}$, and $\lambda$ is the scalar representation of~$\mathcal{L}$ (for finite-dimensional problems, $\lambda$ is a chosen positive scalar; for differential operators, $\lambda$ incorporates the discretization).

The series converges if $|1 + \sigma| < 1$, i.e., $0 < \mathcal{N}'_0/\lambda < 2$ for $c_0 = -1$.

\begin{remark}
With $\lambda = \mathcal{N}'_0$ (Newton): $\sigma = -1$, hence $z_2 = 0$ and $z_3$ reduces to the Halley correction. The formulas~\eqref{eq:z123} generalize Newton--Halley to arbitrary~$\lambda$.
\end{remark}


%% ======================================================================
\section{Level~1: RBF identification by homotopy}
\label{sec:identification}
%% ======================================================================

\subsection{Problem statement}

Given data $\{(y_i, f_i)\}_{i=1}^N$ (where $f_i$ is measured or estimated from $\dot{y}_i + f_i = u_i$), find $\boldsymbol{\theta} = (\mathbf{w}, \mathbf{c}, \boldsymbol{\lambda})$ such that:
\begin{equation}
\mathcal{N}_{\mathrm{id}}(\boldsymbol{\theta}) = \boldsymbol{\Phi}(\mathbf{c}, \boldsymbol{\lambda})\,\mathbf{w} - \mathbf{f} = \mathbf{0},
\label{eq:Nid}
\end{equation}
where $\Phi_{ij} = \varphi(\|y_i - c_j\|/\lambda_j)$ is the RBF design matrix.

\subsection{Bilinear structure}

The residual~\eqref{eq:Nid} is \emph{linear} in~$\mathbf{w}$ and \emph{nonlinear} in~$(\mathbf{c}, \boldsymbol{\lambda})$. This bilinear structure allows a two-stage approach.

\begin{proposition}[Linear subproblem]\label{prop:linear}
For fixed $(\mathbf{c}, \boldsymbol{\lambda})$, the optimal weights are:
\begin{equation}
\mathbf{w}^* = (\boldsymbol{\Phi}^T\boldsymbol{\Phi})^{-1}\boldsymbol{\Phi}^T\mathbf{f}.
\label{eq:ls}
\end{equation}
This is the homotopy correction $z_1$ with $\mathcal{L} = I$ (identity) applied to the linear subproblem. Since $\partial^2\mathcal{N}/\partial w^2 = 0$, the corrections $z_2 = z_3 = 0$ and the solution is exact in one step.
\end{proposition}

\subsection{Nonlinear subproblem: centers}

With $\mathbf{w} = \mathbf{w}^*(\mathbf{c}, \boldsymbol{\lambda})$ substituted, the residual becomes a function of $(\mathbf{c}, \boldsymbol{\lambda})$ only. The Jacobian with respect to center~$c_j$ is:
\begin{equation}
\frac{\partial\mathcal{N}_i}{\partial c_j} = w_j\,\varphi'\!\left(\frac{\|y_i - c_j\|}{\lambda_j}\right)\,\frac{c_j - y_i}{\lambda_j\,\|y_i - c_j\|},
\label{eq:jac-c}
\end{equation}
and the Hessian:
\begin{equation}
\frac{\partial^2\mathcal{N}_i}{\partial c_j^2} = w_j\left[\varphi''\!\left(\frac{r_{ij}}{\lambda_j}\right)\frac{1}{\lambda_j^2} - \varphi'\!\left(\frac{r_{ij}}{\lambda_j}\right)\frac{1}{\lambda_j\,r_{ij}}\right],
\label{eq:hess-c}
\end{equation}
where $r_{ij} = \|y_i - c_j\|$. The homotopy corrections on the centers are:

\textbf{$z_1$ (Newton):}
\begin{equation}
\Delta\mathbf{c}_1 = -(J_c^T J_c)^{-1}\,J_c^T\,\mathcal{N}_0,
\label{eq:z1-c}
\end{equation}
where $J_c$ is the Jacobian matrix with entries~\eqref{eq:jac-c}.

\textbf{$z_2$ (Halley):} After updating $\mathbf{c} \leftarrow \mathbf{c} + \Delta\mathbf{c}_1$ and recomputing $\mathbf{w}^*$ and the residual~$\mathcal{N}_1$:
\begin{equation}
\Delta c_{j,2} = -\frac{1}{2}\,\frac{(\mathcal{N}_1^T e_j)^2\,H_{jj}}{(J_c^T e_j)^3},
\label{eq:z2-c}
\end{equation}
where $H_{jj}$ is the diagonal Hessian element~\eqref{eq:hess-c} summed over data points, and $e_j$ is the $j$-th unit vector in parameter space.

\subsection{Nonlinear subproblem: widths}

The Jacobian with respect to width~$\lambda_j$ is:
\begin{equation}
\frac{\partial\mathcal{N}_i}{\partial\lambda_j} = -w_j\,\varphi'\!\left(\frac{r_{ij}}{\lambda_j}\right)\,\frac{r_{ij}}{\lambda_j^2}.
\label{eq:jac-lambda}
\end{equation}

The correction is:
\begin{equation}
\Delta\boldsymbol{\lambda}_1 = -(J_\lambda^T J_\lambda)^{-1}\,J_\lambda^T\,\mathcal{N}_1.
\label{eq:z1-lambda}
\end{equation}

\subsection{Initialization by K-means}

The initial centers $\mathbf{c}^{(0)}$ are obtained by K-means clustering on $\{y_i\}$, and the initial widths by:
\begin{equation}
\lambda_j^{(0)} = \alpha\,\min_{l \neq j}\|c_j^{(0)} - c_l^{(0)}\|, \qquad \alpha \in [0.5, 2].
\label{eq:init-lambda}
\end{equation}

This provides a reasonable starting point from which the homotopy corrections converge.

\subsection{Complete identification algorithm}

\begin{algorithm}[H]
\caption{HAM-RBF Identification}
\label{alg:identification}
\begin{algorithmic}[1]
\REQUIRE Data $\{(y_i, f_i)\}_{i=1}^N$, number of centers $M$, tolerance $\varepsilon$
\STATE $\mathbf{c}^{(0)} \leftarrow$ K-means$(y_1,\ldots,y_N; M)$
\STATE $\boldsymbol{\lambda}^{(0)} \leftarrow$ nearest-neighbor distances $\times\,\alpha$
\FOR{$k = 0, 1, 2, \ldots$}
    \STATE \textbf{Step~1 (weights, linear, exact):}
    \STATE \quad Build $\boldsymbol{\Phi}^{(k)}$ from $(\mathbf{c}^{(k)}, \boldsymbol{\lambda}^{(k)})$
    \STATE \quad $\mathbf{w}^{(k)} \leftarrow (\boldsymbol{\Phi}^T\boldsymbol{\Phi})^{-1}\boldsymbol{\Phi}^T\mathbf{f}$ \hfill [$z_1$ on $\mathbf{w}$, exact]
    \STATE \quad $\mathcal{N}_0 \leftarrow \boldsymbol{\Phi}^{(k)}\mathbf{w}^{(k)} - \mathbf{f}$
    \IF{$\|\mathcal{N}_0\| < \varepsilon$}
        \STATE \textbf{return} $(\mathbf{w}^{(k)}, \mathbf{c}^{(k)}, \boldsymbol{\lambda}^{(k)})$
    \ENDIF
    \STATE \textbf{Step~2 (centers, Newton $z_1$):}
    \STATE \quad Compute $J_c$ via~\eqref{eq:jac-c}
    \STATE \quad $\Delta\mathbf{c}_1 \leftarrow -(J_c^T J_c)^{-1}J_c^T\mathcal{N}_0$
    \STATE \quad $\mathbf{c}^{(*)} \leftarrow \mathbf{c}^{(k)} + \Delta\mathbf{c}_1$
    \STATE \textbf{Step~3 (centers, Halley $z_2$):}
    \STATE \quad Recompute $\mathbf{w}^{(*)}$, $\mathcal{N}_1$ with $\mathbf{c}^{(*)}$
    \STATE \quad Compute $H_c$ via~\eqref{eq:hess-c}
    \STATE \quad $\Delta\mathbf{c}_2 \leftarrow$ Halley correction~\eqref{eq:z2-c}
    \STATE \quad $\mathbf{c}^{(**)} \leftarrow \mathbf{c}^{(*)} + \Delta\mathbf{c}_2$
    \STATE \textbf{Step~4 (widths, Newton $z_1$):}
    \STATE \quad Recompute $\mathbf{w}^{(**)}$, $\mathcal{N}_2$ with $\mathbf{c}^{(**)}$
    \STATE \quad Compute $J_\lambda$ via~\eqref{eq:jac-lambda}
    \STATE \quad $\Delta\boldsymbol{\lambda}_1 \leftarrow -(J_\lambda^T J_\lambda)^{-1}J_\lambda^T\mathcal{N}_2$
    \STATE \quad $\boldsymbol{\lambda}^{(k+1)} \leftarrow \boldsymbol{\lambda}^{(k)} + \Delta\boldsymbol{\lambda}_1$
    \STATE \quad $\mathbf{c}^{(k+1)} \leftarrow \mathbf{c}^{(**)}$
\ENDFOR
\end{algorithmic}
\end{algorithm}

\begin{remark}
At each outer iteration, the weights $\mathbf{w}$ are recomputed exactly (linear solve, cost $O(NM^2)$). The homotopy corrections are applied only to the $2M$ nonlinear parameters, not to all $3M$. This is not Variable Projection~\cite{GolubPereyra1973} with gradient descent on the nonlinear part---it is Variable Projection with \emph{homotopy corrections} (Newton--Halley) that exploit curvature information through $z_2$.
\end{remark}


%% ======================================================================
\section{Level~2: ODE simulation by homotopy}
\label{sec:simulation}
%% ======================================================================

\subsection{Linear starting solution}

Once the RBF $\hat{f}(y; \boldsymbol{\theta})$ is identified, the ODE becomes:
\begin{equation}
\dot{y} + \hat{f}(y) = u(t), \qquad y(0) = y_0.
\label{eq:ode-rbf}
\end{equation}

We choose the auxiliary linear operator as the linearization of the RBF around a nominal point~$\bar{y}$:
\begin{equation}
\mathcal{L}[y] = \dot{y} + k\,y, \qquad k = \hat{f}'(\bar{y}) = \sum_{j=1}^{M} w_j\,\varphi'\!\left(\frac{\|\bar{y} - c_j\|}{\lambda_j}\right)\frac{1}{\lambda_j}.
\label{eq:L-linear}
\end{equation}

The solution of $\mathcal{L}[y] = u$ with $y(0) = y_0$ is known analytically:
\begin{equation}
y_0(t) = y_0\,e^{-kt} + \int_0^t e^{-k(t-s)}\,u(s)\,ds.
\label{eq:y0-linear}
\end{equation}

This is the Green's function response with exponential kernel $K(t,s) = e^{-k(t-s)}$.

\subsection{Reduced nonlinearity}

The ODE~\eqref{eq:ode-rbf} is rewritten as:
\begin{equation}
\underbrace{\dot{y} + k\,y}_{\mathcal{L}[y]} + \underbrace{\hat{f}(y) - k\,y}_{\tilde{\mathcal{N}}[y]} = u(t).
\label{eq:split}
\end{equation}

The reduced nonlinearity $\tilde{f}(y) = \hat{f}(y) - k\,y$ satisfies $\tilde{f}(\bar{y}) \approx 0$ and $\tilde{f}'(\bar{y}) \approx 0$ by construction. The homotopy deforms from the linear response~$y_0(t)$ to the full nonlinear solution.

\subsection{Integral formulation (derivative-free)}

Following~\cite{Rodrigo2026integral}, we reformulate the ODE as a Volterra integral equation:
\begin{equation}
y(t) = y_0\,e^{-kt} + \int_0^t e^{-k(t-s)}\,[u(s) - \tilde{f}(y(s))]\,ds.
\label{eq:volterra-rbf}
\end{equation}

Discretizing with the trapezoidal rule at step~$k_{\mathrm{step}}$, the integral residual is:
\begin{equation}
g_{\mathrm{sim}}(y_{k_{\mathrm{step}}}) = y_{k_{\mathrm{step}}} + T\,w_{k_{\mathrm{step}}}\,\tilde{f}(y_{k_{\mathrm{step}}}) - T\,w_{k_{\mathrm{step}}}\,u_{k_{\mathrm{step}}} - e^{-kT}\,S_{k_{\mathrm{step}}-1} - y_0\,e^{-k\,t_{k_{\mathrm{step}}}},
\label{eq:g-sim}
\end{equation}
where $S_{k_{\mathrm{step}}}$ is updated via IIR filter:
\begin{equation}
S_{k_{\mathrm{step}}} = e^{-kT}\,S_{k_{\mathrm{step}}-1} + T\,w_{k_{\mathrm{step}}}\,[u_{k_{\mathrm{step}}} - \tilde{f}(y_{k_{\mathrm{step}}})].
\label{eq:iir-sim}
\end{equation}

The Jacobian is:
\begin{equation}
g'_{\mathrm{sim}} = 1 + T\,w_{k_{\mathrm{step}}}\,\tilde{f}'(y_{k_{\mathrm{step}}}) = O(1).
\label{eq:jac-sim}
\end{equation}

\textbf{No division by $T$ appears.} All coefficients are $O(1)$, ensuring numerical stability and noise robustness.

\subsection{Corrections}

The homotopy corrections on $y_{k_{\mathrm{step}}}$ use~\eqref{eq:z123} with:
\begin{equation}
\mathcal{N}_0 = g_{\mathrm{sim}}, \quad \mathcal{N}'_0 = g'_{\mathrm{sim}}, \quad \mathcal{N}''_0 = T\,w_{k_{\mathrm{step}}}\,\tilde{f}''(y_{k_{\mathrm{step}}}), \quad \lambda = g'_{\mathrm{sim}}.
\end{equation}

With $\lambda = g'_{\mathrm{sim}}$ (Newton): $\sigma = -1$, $z_2 = 0$, and $z_3$ is the Halley correction. Three evaluations per step. Complexity: $O(M)$ per step (one RBF evaluation).

\subsection{Derivatives of the RBF}

The derivatives of $\tilde{f}(y) = \hat{f}(y) - ky$ needed for the corrections are:
\begin{align}
\tilde{f}'(y) &= \sum_{j=1}^{M} w_j\,\varphi'\!\left(\frac{|y - c_j|}{\lambda_j}\right)\frac{\mathrm{sgn}(y - c_j)}{\lambda_j} - k, \label{eq:df-rbf} \\
\tilde{f}''(y) &= \sum_{j=1}^{M} w_j\,\varphi''\!\left(\frac{|y - c_j|}{\lambda_j}\right)\frac{1}{\lambda_j^2}, \label{eq:d2f-rbf} \\
\tilde{f}'''(y) &= \sum_{j=1}^{M} w_j\,\varphi'''\!\left(\frac{|y - c_j|}{\lambda_j}\right)\frac{\mathrm{sgn}(y - c_j)}{\lambda_j^3}. \label{eq:d3f-rbf}
\end{align}

For Gaussian RBFs, $\varphi(r) = e^{-r^2}$: $\varphi'(r) = -2r\,e^{-r^2}$, $\varphi''(r) = (4r^2 - 2)\,e^{-r^2}$, etc. All derivatives are closed-form and smooth---no finite differences needed.


%% ======================================================================
\section{Level~3: Online refinement by homotopy}
\label{sec:online}
%% ======================================================================

\subsection{Problem}

During operation, new data $(t_{\mathrm{new}}, y_{\mathrm{new}}, u_{\mathrm{new}})$ arrive. The current RBF model $\hat{f}$ may be inaccurate in regions not covered by the training data. The tracking error:
\begin{equation}
\mathcal{N}_{\mathrm{on}}(\boldsymbol{\theta}) = \dot{y}_{\mathrm{new}} + \hat{f}(y_{\mathrm{new}}; \boldsymbol{\theta}) - u_{\mathrm{new}}
\label{eq:N-online}
\end{equation}
should be driven to zero by updating~$\boldsymbol{\theta}$.

\subsection{Online corrections}

The same three-step structure applies:

\textbf{Step~1 (weights):} With new data appended, re-solve the linear system. If the design matrix is maintained in factored form ($QR$ or Cholesky), the update is $O(NM)$ via rank-one update.

\textbf{Step~2 (centers):} Apply $z_1$ (Newton) on the centers using the accumulated tracking error. Only centers near the new data region are updated (local correction).

\textbf{Step~3 (widths):} Apply $z_1$ on the widths.

\begin{remark}
The online level uses the same formulas as the identification level but operates on a \emph{sliding window} of recent data. The homotopy corrections are incremental---they refine the existing model, not rebuild it from scratch.
\end{remark}


%% ======================================================================
\section{Unification: three levels, one equation}
\label{sec:unification}
%% ======================================================================

\begin{table}[H]
\centering
\caption{Three levels of homotopy applied to RBF-based nonlinear systems.}
\label{tab:levels}
\rowcolors{2}{lightgray}{white}
\begin{tabular}{C{0.8cm} L{2.5cm} C{1.3cm} L{3.2cm} L{2.4cm} L{2.0cm}}
\toprule
\textbf{Level} & \textbf{Problem} & $\boldsymbol{\mathcal{L}}$ & $\boldsymbol{\mathcal{N}}$ & $\boldsymbol{\phi}$ & $\boldsymbol{\lambda}$ \\
\midrule
1a & Weights & $I$ & $\boldsymbol{\Phi}\mathbf{w} - \mathbf{f}$ & $\mathbf{w} \in \mathbb{R}^M$ & $\boldsymbol{\Phi}^T\boldsymbol{\Phi}$ \\
1b & Centers & $I$ & $\mathcal{N}_{\mathrm{id}}(\mathbf{c})$ & $\mathbf{c} \in \mathbb{R}^M$ & $J_c^T J_c$ \\
1c & Widths & $I$ & $\mathcal{N}_{\mathrm{id}}(\boldsymbol{\lambda})$ & $\boldsymbol{\lambda} \in \mathbb{R}^M$ & $J_\lambda^T J_\lambda$ \\
2 & ODE simulation & $d/dt + k$ & $\dot{y} + \tilde{f}(y) - u$ & $y_k \in \mathbb{R}$ & $g'_{\mathrm{sim}}$ \\
3 & Online update & $I$ & tracking error & $(\mathbf{c}, \boldsymbol{\lambda})$ & $J^T J$ \\
\bottomrule
\end{tabular}
\end{table}

\begin{theorem}[Unified homotopy framework]\label{thm:unified}
The five subproblems in Table~\ref{tab:levels} are all instances of the deformation equation~\eqref{eq:deformation}, with corrections given by~\eqref{eq:z123}. Specifically:
\begin{enumerate}[nosep]
\item[\emph{(i)}] Level~1a (weights) is solved exactly in one step ($z_2 = z_3 = 0$).
\item[\emph{(ii)}] Levels~1b, 1c, and~3 converge if $0 < \mathcal{N}'_0/\lambda < 2$ at each iteration.
\item[\emph{(iii)}] Level~2 converges if $\tilde{f}$ is Lipschitz and $T$ is small enough.
\item[\emph{(iv)}] No gradients, no learning rates, and no backpropagation are used at any level.
\item[\emph{(v)}] The same code implementing~\eqref{eq:z123} serves all five subproblems.
\end{enumerate}
\end{theorem}

\begin{proof}[Proof sketch]
(i)~follows from $\partial^2\mathcal{N}/\partial w^2 = 0$ (linearity in weights). (ii)~is the standard HAM convergence condition applied to each subproblem. (iii)~follows from the integral formulation: the Volterra operator is compact, and the Neumann series converges for finite~$T_f$ by quasi-nilpotency~\cite{Kress2014}. (iv)~is by construction: all corrections are computed from $\mathcal{N}_0$, $\mathcal{N}'_0$, $\mathcal{N}''_0$ analytically. (v)~is the structural observation that~\eqref{eq:z123} depends only on $(\gamma, \sigma, \mathcal{N}''_0)$, which are scalars regardless of the dimension of~$\phi$.
\end{proof}


%% ======================================================================
\section{The role of the linear starting solution}
\label{sec:linear-start}
%% ======================================================================

The quality of the homotopy depends critically on the initial approximation. For Level~2 (simulation), the linear starting solution
\begin{equation}
y_0(t) = y_0\,e^{-kt} + \int_0^t e^{-k(t-s)}\,u(s)\,ds
\label{eq:y0-repeat}
\end{equation}
has three desirable properties:

\textbf{(a)~It satisfies the initial condition exactly:} $y_0(0) = y_0$.

\textbf{(b)~It captures the dominant dynamics:} if $\hat{f}(y) \approx k\,y$ near the operating point, then $y_0(t) \approx y(t)$ and the corrections $z_1, z_2, z_3$ are small.

\textbf{(c)~It defines the kernel:} $K(t,s) = e^{-k(t-s)}$ is the Green's function of $\mathcal{L} = d/dt + k$. The kernel encodes the linear time constant $\tau = 1/k$ of the system.

The optimal choice is $k = \hat{f}'(y_{\mathrm{eq}})$, the linearization around the equilibrium. This minimizes $\|\tilde{f}\| = \|\hat{f} - k\,\mathrm{id}\|$ near the operating point, making the corrections as small as possible.

For systems with multiple operating regimes, $k$ can be adapted: use the linearization of the RBF at the current state $y_{k_{\mathrm{step}}-1}$. This makes $\lambda = g'_{\mathrm{sim}} \approx 1$ (well-conditioned), ensuring rapid convergence.


%% ======================================================================
\section{Connection to existing methods}
\label{sec:connections}
%% ======================================================================

\begin{table}[H]
\centering
\caption{The unified framework recovers classical methods as special cases.}
\label{tab:connections}
\rowcolors{2}{lightgray}{white}
\begin{tabular}{L{4cm} L{3.5cm} L{5cm}}
\toprule
\textbf{Choice} & \textbf{Method recovered} & \textbf{Reference} \\
\midrule
$\mathcal{L} = I$, $z_1$ only on $\mathbf{w}$ & Least squares & Golub--Pereyra~\cite{GolubPereyra1973} \\
$\mathcal{L} = I$, $z_1$ on $(\mathbf{c}, \boldsymbol{\lambda})$ & Gauss--Newton & Bj\"orck~\cite{Bjorck1996} \\
$\mathcal{L} = I$, $z_1 + z_2$ on $(\mathbf{c}, \boldsymbol{\lambda})$ & Halley's method & Guti\'errez--Hern\'andez~\cite{Gutierrez1997} \\
$\mathcal{L} = d/dt$, integral residual & Picard iteration & Atkinson~\cite{Atkinson1997} \\
$\mathcal{L} = d/dt + k$, $z_1$ only & Newton on integral ODE & This paper \\
$\mathcal{L} = d/dt + k$, $z_1 + z_2 + z_3$ & Newton--Halley on integral ODE & This paper \\
Full framework (Levels 1--3) & \textbf{Unified HAM-RBF} & \textbf{This paper} \\
\bottomrule
\end{tabular}
\end{table}

The contribution is not any single row of Table~\ref{tab:connections} but the observation that \emph{all rows are instances of the same equation}~\eqref{eq:deformation}.


%% ======================================================================
\section{Computational complexity}
\label{sec:complexity}
%% ======================================================================

\begin{table}[H]
\centering
\caption{Computational cost per outer iteration / per time step.}
\label{tab:complexity}
\rowcolors{2}{lightgray}{white}
\begin{tabular}{L{4.5cm} C{2.5cm} L{4.5cm}}
\toprule
\textbf{Operation} & \textbf{Cost} & \textbf{Notes} \\
\midrule
Weights (linear solve) & $O(NM^2)$ & One-time per outer iteration \\
Jacobian $J_c$ & $O(NM)$ & Sparse if RBFs are local \\
Newton $z_1$ on $\mathbf{c}$ & $O(M^3)$ & Normal equations, $M \ll N$ \\
Halley $z_2$ on $\mathbf{c}$ & $O(NM)$ & Diagonal Hessian approximation \\
ODE step ($z_1 + z_2 + z_3$) & $O(M)$ & One RBF evaluation per correction \\
IIR accumulator update & $O(1)$ & One multiply-add \\
\bottomrule
\end{tabular}
\end{table}

The identification phase (Level~1) costs $O(NM^2)$ per outer iteration, dominated by the linear solve. Typically 3--5 outer iterations suffice. The simulation phase (Level~2) costs $O(M)$ per time step---embeddable on microcontrollers.


%% ======================================================================
\section{Convergence analysis}
\label{sec:convergence}
%% ======================================================================

\subsection{Level~1: identification}

\begin{proposition}[Convergence of center correction]\label{prop:conv-centers}
Let $\hat{f}$ be a Gaussian RBF with $M$ centers. If the initial centers $\mathbf{c}^{(0)}$ (from K-means) are within the basin of attraction of the true centers, then the Newton correction~\eqref{eq:z1-c} converges quadratically and the Halley correction~\eqref{eq:z2-c} converges cubically, provided:
\begin{equation}
\|J_c^{-1}\|\,\|H_c\|\,\|\Delta\mathbf{c}_1\| < 2.
\label{eq:conv-cond}
\end{equation}
\end{proposition}

\subsection{Level~2: simulation}

\begin{proposition}[Convergence of ODE step]\label{prop:conv-ode}
For the integral residual~\eqref{eq:g-sim} with $\lambda = g'_{\mathrm{sim}}$ (Newton):
\begin{enumerate}[nosep]
\item[\emph{(i)}] $z_1$ gives quadratic convergence to the fixed point of the discrete equation.
\item[\emph{(ii)}] The discretization error is $O(T^2)$ (trapezoidal).
\item[\emph{(iii)}] The integral residual is $O(1)$, not $O(1/T)$: noise in the data is not amplified.
\end{enumerate}
\end{proposition}

\subsection{Combined convergence}

The identification error in $\hat{f}$ propagates to the simulation as a model error. If $\|\hat{f} - f\|_\infty = \epsilon$, then the simulation error satisfies:
\begin{equation}
\|y - \hat{y}\|_\infty \leq C\,\epsilon\,T_f\,e^{L T_f},
\label{eq:combined-error}
\end{equation}
where $L$ is the Lipschitz constant of $f$ and $T_f$ is the simulation horizon. Reducing $\epsilon$ (better RBF fit) directly improves simulation accuracy.


%% ======================================================================
\section{Discussion}
\label{sec:discussion}
%% ======================================================================

\textbf{Why homotopy and not backpropagation?} Backpropagation computes $\partial\mathcal{N}/\partial\boldsymbol{\theta}$ (the gradient) and takes a step of size~$\eta$ (learning rate) in the negative gradient direction. The homotopy correction $z_1 = -\mathcal{N}_0/\mathcal{N}'_0$ is Newton's method---it uses the \emph{same} Jacobian but divides by it instead of multiplying. Newton converges quadratically; gradient descent converges linearly. The Halley correction $z_2$ adds curvature information, reaching cubic convergence. No learning rate is needed because the step size is determined by the local geometry of~$\mathcal{N}$.

\textbf{Why separate linear and nonlinear parameters?} The RBF residual is bilinear: linear in~$\mathbf{w}$, nonlinear in~$(\mathbf{c}, \boldsymbol{\lambda})$. Solving $\mathbf{w}$ exactly by linear algebra (one step) and applying homotopy only to $(\mathbf{c}, \boldsymbol{\lambda})$ reduces the nonlinear problem from $3M$ to $2M$ unknowns and exploits the bilinear structure optimally. This is the Variable Projection principle~\cite{GolubPereyra1973}, but with Newton--Halley instead of gradient-based optimization on the nonlinear part.

\textbf{Why the integral formulation for simulation?} The BDF residual $g_{\mathrm{BDF}} = (3y_k - 4y_{k-1} + y_{k-2})/(2T) + \hat{f}(y_k) - u_k$ contains $1/T$, amplifying sensor noise. The integral residual~\eqref{eq:g-sim} contains only $T$ as a multiplicative factor, attenuating noise. Both have $O(T^2)$ discretization error, but the noise propagation properties differ fundamentally.

\textbf{Limitations.} The convergence of Level~1 (identification) depends on the K-means initialization being in the basin of attraction of the true centers. For highly multimodal problems, multiple restarts or global optimization may be needed for the initial centers. The convergence of Level~2 (simulation) requires $T$ small enough that the 3-term homotopy series converges; for stiff systems, this may be restrictive.

\textbf{Connection to the Homotopy Analysis Method.} Liao's HAM~\cite{Liao2003} was developed for computing \emph{analytic} series solutions to differential equations. The present work uses the HAM deformation equation as an \emph{algorithmic} framework---the series is truncated at 3 terms and evaluated numerically at each step. The convergence-control parameter $c_0$ and the freedom of choosing~$\mathcal{L}$ survive in the discrete setting as the parameter~$\lambda$ and the choice of kernel~$K$.


%% ======================================================================
\section{Conclusions}
\label{sec:conclusions}
%% ======================================================================

\begin{enumerate}[nosep]
\item Liao's deformation equation $(1-q)\mathcal{L}[\phi - \phi_0] = q\,c_0\,\mathcal{N}[\phi]$ is a universal root-finding framework. Its corrections $z_1$ (Newton), $z_2$ (Halley), $z_3$ (third order) are parametrized by $(\gamma, \sigma)$ and apply to any $\mathcal{N}$.

\item For $\dot{y} + \text{RBF}(y) = u(t)$, three problems---identification, simulation, and online refinement---are solved by the same corrections with different choices of $(\mathcal{L}, \mathcal{N}, \phi)$.

\item The bilinear structure of the RBF is exploited: weights are solved exactly (one linear step), centers and widths are corrected by Newton--Halley (3--5 iterations).

\item The linear starting solution $y_0(t) = y_0 e^{-kt} + \int_0^t e^{-k(t-s)}u(s)\,ds$ provides an analytically known initial approximation for the simulation homotopy, with kernel $K = e^{-k(t-s)}$ encoding the system's linear time constant.

\item The integral formulation for simulation eliminates numerical derivatives, with all residual coefficients at $O(1)$, ensuring noise robustness.

\item No gradients, no learning rates, no backpropagation. All corrections are closed-form.
\end{enumerate}


%% ======================================================================
\begin{thebibliography}{99}
%% ======================================================================

\bibitem{Liao2003}
S.\,J.\ Liao, \textit{Beyond Perturbation: Introduction to the Homotopy Analysis Method}, Chapman \& Hall/CRC, 2003.

\bibitem{Liao2012}
S.\,J.\ Liao, \textit{Homotopy Analysis Method in Nonlinear Differential Equations}, Springer, 2012.

\bibitem{Liao2010}
S.\,J.\ Liao, ``An optimal homotopy-analysis approach for strongly nonlinear differential equations,'' \textit{Commun.\ Nonlinear Sci.\ Numer.\ Simul.}, vol.~15, pp.~2003--2016, 2010.

\bibitem{Haykin2009}
S.\ Haykin, \textit{Neural Networks and Learning Machines}, 3rd ed., Prentice Hall, 2009.

\bibitem{Park1991}
J.\ Park, I.\,W.\ Sandberg, ``Universal approximation using radial-basis-function networks,'' \textit{Neural Comput.}, vol.~3, pp.~246--257, 1991.

\bibitem{Hairer1993}
E.\ Hairer, S.\,P.\ N\o rsett, G.\ Wanner, \textit{Solving Ordinary Differential Equations I: Nonstiff Problems}, 2nd ed., Springer, 1993.

\bibitem{Ljung1999}
L.\ Ljung, \textit{System Identification: Theory for the User}, 2nd ed., Prentice Hall, 1999.

\bibitem{GolubPereyra1973}
G.\,H.\ Golub, V.\ Pereyra, ``The differentiation of pseudo-inverses and nonlinear least squares problems whose variables separate,'' \textit{SIAM J.\ Numer.\ Anal.}, vol.~10, pp.~413--432, 1973.

\bibitem{Bjorck1996}
\AA.\ Bj\"orck, \textit{Numerical Methods for Least Squares Problems}, SIAM, 1996.

\bibitem{Gutierrez1997}
J.\,M.\ Guti\'errez, M.\,A.\ Hern\'andez, ``A family of Chebyshev--Halley type methods in Banach spaces,'' \textit{Bull.\ Aust.\ Math.\ Soc.}, vol.~55, pp.~113--130, 1997.

\bibitem{Atkinson1997}
K.\,E.\ Atkinson, \textit{The Numerical Solution of Integral Equations of the Second Kind}, Cambridge University Press, 1997.

\bibitem{Kress2014}
R.\ Kress, \textit{Linear Integral Equations}, 3rd ed., Springer, 2014.

\bibitem{Rodrigo2026integral}
R.\,H.\ Rodrigo, H.\,D.\ Pati\~no, S.\,J.\ Liao, ``Integral homotopy regressors: unifying continuous and discrete HAM via compact Volterra operators,'' submitted to \textit{Commun.\ Nonlinear Sci.\ Numer.\ Simul.}, 2026.

\bibitem{Rodrigo2024}
R.\,H.\ Rodrigo, H.\,D.\ Pati\~no, ``Homotopy-based functional neural network (HFNN) for nonlinear system identification,'' submitted to \textit{Neurocomputing}, 2025.

\bibitem{Chen1991}
S.\ Chen, C.\,F.\,N.\ Cowan, P.\,M.\ Grant, ``Orthogonal least squares learning algorithm for radial basis function networks,'' \textit{IEEE Trans.\ Neural Netw.}, vol.~2, pp.~302--309, 1991.

\bibitem{Buhmann2003}
M.\,D.\ Buhmann, \textit{Radial Basis Functions: Theory and Implementations}, Cambridge University Press, 2003.

\bibitem{Poggio1990}
T.\ Poggio, F.\ Girosi, ``Networks for approximation and learning,'' \textit{Proc.\ IEEE}, vol.~78, pp.~1481--1497, 1990.

\bibitem{Broomhead1988}
D.\,S.\ Broomhead, D.\ Lowe, ``Multivariable functional interpolation and adaptive networks,'' \textit{Complex Syst.}, vol.~2, pp.~321--355, 1988.

\bibitem{Nelles2001}
O.\ Nelles, \textit{Nonlinear System Identification}, Springer, 2001.

\bibitem{Sjoberg1995}
J.\ Sj\"oberg, Q.\ Zhang, L.\ Ljung, \textit{et al.}, ``Nonlinear black-box modeling in system identification: a unified overview,'' \textit{Automatica}, vol.~31, pp.~1691--1724, 1995.

\end{thebibliography}

\end{document}
